% Source: source_md/intelligence_part2_combined_ja.md
\section{非閉包定理 — 推論層 $L_F$ に内部不動点が存在しないこと}


本章は本論文の第一中心である(第二中心は \S{}5)。VEDの基本構造の上で、推論層 $L_F$ に\textbf{内部不動点が存在しないこと}を構造的命題として導出し、その動力学的帰結を述べる。

この帰結は、知性が推論層 $L_F$ 単独で閉じることを構造的に禁ずる。すなわち、停止の外部化は設計上の選択ではなく、推論の構造そのものから導かれる必然である。\S{}5 では、この「単独では閉じない」が「ログ層 $L_R$ と創発停止層 $L_E$ を通じてなら閉じうる」へと精密化される。本章はその前提となる、$L_F$ 単独の限界を確立する。

\begin{figure}[H]
\centering
\includegraphics[width=\linewidth]{fig1_c_landscape_ja.pdf}
\caption{C-地景と推論軌道。推論は閉包度の地形の上を勾配 $\nabla C$ に沿って動き続けるが、構造的にどこにも止まれない。}
\label{fig:part2-c-landscape}
\figurenote{横軸は状態空間、縦軸は閉包度 $C$ である。青い軌道は推論演算子 $F$ の反復を、赤い境界は到達不能な $C_{\max}$ を示す。}{固定された $\theta$ の下での領域内移動 (A) では、$L_F$ は内部不動点を持たない。本図は非閉包定理の幾何学的表現であり、後に $\mu_\theta$ による領域組み替え (B) と $L_E$ の創発が必要になる理由を示す。}
\end{figure}

\bigskip


\subsection{推論演算子の前提と正当化測度}


\S{}2で導入した推論演算子 $F$ を、本章での議論に必要な範囲で再定式化する。

\textbf{前提 (P1) — 正当化拡張性}:
$F$ は状態 $x_n$ から、より正当化された状態 $x_{n+1}$ への遷移を生成する。

\[
x_{n+1} = F(x_n)
\]

ここで $x_n$ は仮説・解釈・理解のいずれかの状態を表す。$F$ の具体的実装(演繹的・確率的・統計的)は本議論では問わない。問うのは、$F$ が\textbf{正当化を拡張する操作}であるという最小性質のみである。

\textbf{前提 (P2) — 正当化測度の存在}:
ある測度 $J: \mathcal{X} \to \mathbb{R}$ が存在し、$F$ の作用によって単調に増大する。すなわち任意の許容状態 $x \in \mathcal{X}$ について:

\[
J(F(x)) > J(x)
\]

この測度 $J$ は、推論連鎖の進行を測る量であって、真理性を測る量ではない。後段(\S{}10)で明確になるように、ゲティア構造の存在により、$J$ と真理は同一視できない。

\textbf{前提 (P3) — VED接地}:
推論演算子 $F$ は、VEDの流れ場 $J_C = -\xi \nabla C$ の知性層実現として解釈される。すなわち、推論は閉包度 $C$ の勾配に沿った状態遷移である。さらに \S{}2.1 の通り、$F$ は Bio-IFGT の状態発展演算子 $F(x;\theta)$ の知性層特殊形であり、運動可能領域の拘束パラメータ $\theta$ の下で作用する。本章の議論は固定された $\theta$ の下、すなわち単一の運動可能領域 $\Omega_\theta$ の内部における $F$ の挙動を扱う(\S{}2.7 の区別における (A) 領域内移動の範囲)。$\theta$ 自体の組み替え (B) は本章の射程外であり、\S{}5 で扱われる。

これら三前提のもとで主定理を述べる。

\bigskip


\subsection{主定理 — 内部不動点の不在}


\textbf{定理 3.1 (非閉包定理 / Non-Closure Theorem)}

\begin{quote}
前提 (P1), (P2) を満たす任意の推論演算子 $F$ について、内部不動点 $x^*$ は存在しない。すなわち:

$$
\forall x \in \mathcal{X}, \quad F(x) \neq x
$$
\end{quote}


ここで「内部不動点」とは、$F$ の作用域内で $F(x^*) = x^*$ を満たす状態を指す。

この定理は、推論が\textbf{自己内で停止できないこと}を構造的に表明する。

\bigskip


\subsection{証明スケッチ}


仮に $F(x^*) = x^*$ を満たす不動点 $x^* \in \mathcal{X}$ が存在すると仮定する。

前提 (P2) により、任意の許容状態について:

\[
J(F(x)) > J(x)
\]

特に $x = x^*$ について:

\[
J(F(x^*)) > J(x^*)
\]

しかし不動点の仮定により $F(x^*) = x^*$、したがって:

\[
J(x^*) > J(x^*)
\]

これは矛盾。ゆえに不動点 $x^*$ は存在しない。 $\blacksquare$

証明は短く、本質的には「正当化拡張性」という $F$ の最小性質のみに依拠する。すなわち、推論を\textbf{正当化を進める操作}と定義した時点で、内部不動点は構造的に排除されている。

\bigskip


\subsection{失敗ではなく構造的性質}


この結果はしばしば「推論の限界」と誤解される。本論文の立場では、これは限界ではなく\textbf{設計}である。

推論演算子 $F$ は「動くように」設計された操作である。停止する設計ではない。これは欠陥ではなく、推論という操作の定義そのものに由来する性質である。

類比として、流体の連続方程式 $\partial \rho / \partial t + \nabla \cdot (\rho \mathbf{v}) = 0$ が「停止」しないことは、流体の欠陥ではない。流体は流れる対象として定義されている。同様に、推論は進行する操作として定義されている。

VEDの語彙で述べれば: 推論演算子 $F$ は閉包度勾配 $\nabla C$ に沿った流れであり、流れである限り静止しない。静止は流れの完成ではなく終了である。

\bigskip


\subsection{内部停止述語の不可能性}


定理 3.1 から、次の系が直ちに従う。

\textbf{系 3.2 (内部停止述語の不可能性)}

\begin{quote}
推論内部に停止述語 $S: \mathcal{X} \to \{0, 1\}$ を定義することはできない。任意のそのような述語の定義は、無限後退を生成する。
\end{quote}


理由は次の通り。仮に内部停止述語 $S$ を定義すると、$S(x) = 1$ となる状態 $x$ は不動点的性質を持たねばならない(これ以上推論を進める必要がないとの判断状態)。しかし定理 3.1 により、$F$ の作用域内にそのような状態は存在しない。

$S$ を「停止判断のための補助推論」として定義する場合、$S$ 自体が推論演算子 $F'$ として振る舞い、$F'$ にも定理 3.1 が適用される。すなわち $F'$ にも内部不動点はなく、$S$ を確定するためにはさらに $S'$ が必要となる。この連鎖は終わらない。

これが\textbf{正当化の無限後退} (justificatory regress) のVED的定式化である。古来の認識論的問題は、構造的事実の表現であった。

\bigskip


\subsection{VED対応 — $C_{max}$ 境界との同型}


\textbf{系 3.3 ($C_{max}$ 不到達性との同型)}

\begin{quote}
推論演算子 $F$ の内部不動点不在は、VED基本式における $C_{max}$ 不到達性と構造的に同型である。
\end{quote}


VED基本式:

\[
\frac{\partial C}{\partial \tau} + \nabla \cdot J_C = a\Delta\left(1 - \frac{C}{C_*}\right) - \Gamma C - \frac{\kappa_B}{C_{max} - C}
\]

の最終項 $-\kappa_B/(C_{max} - C)$ は、$C \to C_{max}$ で発散する。これは完全閉包への到達が動力学的に不可能であることを示す。

知性層の対応:

\begin{longtable}{@{}p{0.30\linewidth}p{0.30\linewidth}@{}}
\toprule
VED物理層 & 知性層 \\
\midrule
\endfirsthead
\toprule
VED物理層 & 知性層 \\
\midrule
\endhead
$C \to C_{max}$ の発散 & $F$ の内部不動点不在 \\
$\kappa_B/(C_{max}-C)$ 障壁 & 正当化の無限後退 \\
準閉包 $\Omega_\theta$ & 操作可能推論領域 \\
差分の地平線 (Vol.4) & 内部停止の不可能性 \\
\bottomrule
\end{longtable}


両者は別領域の現象だが、構造的には同一の動力学的命題を表現している。これは Part I で示した領域横断ホモロジーの一例である。

\bigskip


\subsection{動力学的解釈と分岐爆発との接続}


定理 3.1 は、推論の動力学を相空間で記述すると、より直観的に見える。

推論演算子 $F$ は状態空間 $\mathcal{X}$ 上の流れ場を生成する。$J$ が単調増大する以上、すべての軌道は $J$ の増大方向に進む。固定点を持たない流れは、$\mathcal{X}$ 内で次の三つの挙動のいずれかを示す:

\begin{enumerate}
\item \textbf{無限後退} — 軌道が境界に向かって発散する
\item \textbf{周期軌道} — 軌道が閉曲線上を巡回する(ただし $J$ 単調性により排除される)
\item \textbf{分岐} — 軌道が複数の方向に分岐し、各分岐がさらに分岐する
\end{enumerate}


実際の推論動力学では、(1)と(3)が結合する。すなわち、推論は\textbf{発散しながら分岐する}。

これは Part I で導入した\textbf{分岐爆発} (branching explosion) と同一現象である。Part I では概念分岐の指数的増殖として記述されたが、\S{}3 の文脈では「内部停止条件の不在が分岐を停止できない」という形で再記述される。

推論能力の増大は、分岐爆発の抑制ではなく、その加速を意味する可能性がある。したがって知性の問題は、推論能力の増強そのものではなく、増大した分岐をいかに停止可能な形で運用するかという問題へ移る。

分岐爆発と内部不動点不在は、同じ構造的事実の二つの側面である:
\begin{itemize}
\item 内部不動点不在 → 軌道は静止しない
\item 分岐爆発 → 軌道は単一でない
\end{itemize}


両者が結合すると、推論は「静止せず、かつ単一でもない」という二重に不安定な動力学を示す。

\bigskip


\subsection{系 — 知性が個体内で閉じることの不可能性}


定理 3.1 と系 3.2 から、本論文の中心的帰結が直接導かれる。

\textbf{系 3.4 (個体閉包不可能性)}

\begin{quote}
いかなる個体システムも、推論層 $L_F$ 単独で推論を閉じることはできない。
\end{quote}


ここで「閉じる」とは、外部参照なしに最終的判断状態に到達することを意味する。系 3.4 は、これが構造的に不可能であることを示す。なお「単独で」という限定が重要である。本系は $L_F$ のみによる閉包を否定するのであって、$L_R$ と $L_E$ を含む三層動力学全体による停止を否定するものではない(\S{}5)。

帰結:

\begin{itemize}
\item 個体ノードが推論の全責任を $L_F$ 内部で担うことを試みれば、必然的に以下のいずれかに陥る:
\item \textbf{過剰確信} (premature closure) — 早すぎる停止
\item \textbf{麻痺} (paralysis) — 無限懐疑
\item \textbf{妄想的閉包} (delusional closure) — 自己封印的合理化
\end{itemize}


これら三つは「弱い知性」の症状ではなく、「停止を外部化していない強い推論」、すなわち $L_E$ を欠いた $L_F$ の必然的引力子である。$L_F$ だけが暴走し、$L_E$ がそれを停止できないとき、推論はこれら三つの引力子のいずれかに落ちる。(なお「病的」「過剰確信」等の語は臨床的診断概念ではなく、停止構造が極限へ落ち込む動力学的状態を指すために拝借したものである。この用語上の断りは \S{}8.5 で詳述する。)

この観察は \S{}11 (不安定領域)で、$L_F$ / $L_R$ / $L_E$ の四重崩壊および $L_E$ 固着として再登場する。高速・高解像・内部停止の知性が必然的に不安定領域を占める理由は、ここで既にその核が証明されている。

\bigskip


\subsection{次章への橋渡し}


定理 3.1 は構造的命題であり、推論の操作的環境を問わない。しかし実際の推論は、応答的環境(他者・実験・社会的反応)と沈黙的環境(自然・過去・未来偶然)の二種類に分けられる。

これら二環境において、内部不動点不在の帰結は異なる形で現れる:

\begin{itemize}
\item 応答的環境では、推論ループは外部応答によって閉じうる(ただし内部不動点ではなく外部停止である)
\item 沈黙的環境では、推論ループは原理的に閉じない
\end{itemize}


次章 \S{}4 では、この区別を C-幾何学の語彙で形式化し、応答チャネルの有無が推論動力学に与える影響を分析する。ただし \S{}2.6 で示した通り、応答的/沈黙的の区別は本体論的範疇ではなく閾値現象であることに注意する。

そこで明らかになるのは、沈黙世界における推論は\textbf{必ず}外部停止演算子を必要とし、その演算子は推論の関数列の外側に位置せざるをえない、ということである。さらに \S{}4 末尾では、無限後退する $L_F$ の運動がログとして堆積し、ログ層 $L_R$ の予兆を生むことが示される。

これらが \S{}5 で展開される三層動力学——ログ層 $L_R$、創発停止層 $L_E$、外部停止演算子 $E$、そして外部参照場 $\Phi_{\text{ext}}$ が個体間 $L_E$ ネットワーク $\Phi_{\text{net}}$ と同型であること——へと繋がる。$L_F$ が単独で閉じないという本章の帰結は、\S{}5 における $L_E$ 創発の必然性の根拠である。

\bigskip


\subsection{要約}


本章で示したのは次の四点である:

\begin{enumerate}
\item \textbf{定理 3.1}: 推論層 $L_F$ に内部不動点が存在しない(正当化拡張性からの構造的帰結)
\item \textbf{系 3.2}: 内部停止述語は無限後退を生む
\item \textbf{系 3.3}: この性質は VED の $C_{max}$ 不到達性と同型である
\item \textbf{系 3.4}: 個体システムは $L_F$ 単独で推論を閉じることができない
\end{enumerate}


これらは「推論層 $L_F$ 単独では構造的に閉じられない」という、本論文の第一中心をなす。本章の証明は固定された運動可能領域 $\Omega_\theta$(固定 $\theta$)の内部、すなわち領域内移動 (A) の範囲で成立する。$\theta$ の組み替え (B) と、ログ層 $L_R$・創発停止層 $L_E$ を通じた停止は、第二中心 \S{}5 で扱われる。

\S{}4 では、停止の外部化がまず応答チャネルの問題として現れることを示し、$L_R$ の予兆を導く。\S{}5 では、$L_F$ が単独で閉じないという本章の帰結を根拠として、創発停止層 $L_E$ がいかにして個体内に外部依存型の停止を成立させるかを展開する。第一中心(閉じない)と第二中心(それでも停止する)が統合されることで、知性の三層動力学が完成する。

\bigskip
